% GENERATED FILE. Edit canonical lessons, then run npm run generate:books.
% canonical-source-hash: fnv1a64:1162c3fe99ae7080
% canonical-lessons: HI-C274-dharkan, HI-C274-saas, HI-C274-khujli, HI-C274-jalan, HI-C274-sujan

\chapter{A heartbeat, Breath, An itch, A burning feeling, A swelling}
\label{ch:hi-a-heartbeat-breath-an-itch-a-burning-feeling-a-swelling}
\hlchaptermodality{274}

\begin{chapteropening}
\textbf{By the end of this chapter:} \emph{I can say a heartbeat, breath, an itch, a burning feeling and a swelling.}

Complete the last lesson of chapter 274: I can say a heartbeat, breath, an itch, a burning feeling and a swelling.
\end{chapteropening}

\section[\texorpdfstring{\hi{धड़कन} (dhaṛkan)}{धड़कन (dhaṛkan)}]{\hi{धड़कन} (dhaṛkan) — a heartbeat}

\label{lesson:HI-C274-dharkan}



\begin{quote}
Before the new one: say the Hindi for a disease, then the Hindi for an illness.
\end{quote}

\subsection*{You'll want to know: \hi{धड़कन}}
\textbf{\hi{धड़कन}} — \emph{dhaṛkan} — \textquotedblleft{}a heartbeat\textquotedblright{}.

\begin{cousinweb}
\textbf{\hi{धड़कन}} (\emph{dhaṛkan}) is a tadbhava, a word inherited and reshaped by sound change. It comes from \emph{*dhaḍ}, \textquotedblleft{}throb, beat\textquotedblright{}, a form CDIAL reconstructs from its descendants (CDIAL 6711). \emph{dhaṛkan} is the noun of \emph{dhaṛaknā}, \textquotedblleft{}to throb\textquotedblright{}.
\end{cousinweb}

\begin{grammarlens}[title={What you've built}]
One more word for this chapter.
\end{grammarlens}

\subsection*{Your turn}
\begin{itemize}
\raggedright
  \item \emph{Say it:} \emph{dhaṛkan}
  \item \emph{Say it:} \emph{dhaṛkan}, once more
  \item \emph{Say it:} say \emph{dhaṛkan}
  \item \emph{Recall:} say the Hindi for a disease, then the Hindi for an illness, then say \emph{dhaṛkan} again
\end{itemize}

\begin{tcolorbox}[breakable,colback=teal!4,colframe=teal!35!black,title={Before you move on}]
What does \hi{धड़कन} mean? (\textquotedblleft{}A heartbeat\textquotedblright{}.) Say it once more.
\end{tcolorbox}
\section[\texorpdfstring{\hi{साँस} (sā̃s)}{साँस (sā̃s)}]{\hi{साँस} (sā̃s) — breath}

\label{lesson:HI-C274-saas}



\begin{quote}
Before the new one: say the Hindi for an illness, then the Hindi for a heartbeat.
\end{quote}

\subsection*{You'll want to know: \hi{साँस}}
\textbf{\hi{साँस}} — \emph{sā̃s} — \textquotedblleft{}breath\textquotedblright{}.

\begin{cousinweb}
\textbf{\hi{साँस}} (\emph{sāṁs}) is a tadbhava, a word inherited from Sanskrit and reshaped by sound change. It comes from \emph{śvāsa}, \textquotedblleft{}breath, breathing\textquotedblright{} (CDIAL 12769).
\end{cousinweb}

\begin{grammarlens}[title={What you've built}]
One more word for this chapter.
\end{grammarlens}

\subsection*{Your turn}
\begin{itemize}
\raggedright
  \item \emph{Say it:} \emph{sā̃s}
  \item \emph{Say it:} \emph{sā̃s}, once more
  \item \emph{Say it:} say \emph{sā̃s}
  \item \emph{Recall:} say the Hindi for an illness, then the Hindi for a heartbeat, then say \emph{sā̃s} again
\end{itemize}

\begin{tcolorbox}[breakable,colback=teal!4,colframe=teal!35!black,title={Before you move on}]
What does \hi{साँस} mean? (\textquotedblleft{}Breath\textquotedblright{}.) Say it once more.
\end{tcolorbox}
\section[\texorpdfstring{\hi{खुजली} (khujlī)}{खुजली (khujlī)}]{\hi{खुजली} (khujlī) — an itch}

\label{lesson:HI-C274-khujli}



\begin{quote}
Before the new one: say the Hindi for a heartbeat, then the Hindi for breath.
\end{quote}

\subsection*{You'll want to know: \hi{खुजली}}
\textbf{\hi{खुजली}} — \emph{khujlī} — \textquotedblleft{}an itch\textquotedblright{}.

\begin{cousinweb}
\textbf{\hi{खुजली}} (\emph{khujlī}) is a tadbhava, a word inherited from Sanskrit and reshaped by sound change. It comes from \emph{kharju}, \textquotedblleft{}itching, scab\textquotedblright{} (CDIAL 3827). CDIAL takes the word from Dravidian. The Sanskrit word is known only from the old word-lists.
\end{cousinweb}

\begin{grammarlens}[title={What you've built}]
One more word for this chapter.
\end{grammarlens}

\subsection*{Your turn}
\begin{itemize}
\raggedright
  \item \emph{Say it:} \emph{khujlī}
  \item \emph{Say it:} \emph{khujlī}, once more
  \item \emph{Say it:} say \emph{khujlī}
  \item \emph{Recall:} say the Hindi for a heartbeat, then the Hindi for breath, then say \emph{khujlī} again
\end{itemize}

\begin{tcolorbox}[breakable,colback=teal!4,colframe=teal!35!black,title={Before you move on}]
What does \hi{खुजली} mean? (\textquotedblleft{}An itch\textquotedblright{}.) Say it once more.
\end{tcolorbox}
\section[\texorpdfstring{\hi{जलन} (jalan)}{जलन (jalan)}]{\hi{जलन} (jalan) — a burning feeling}

\label{lesson:HI-C274-jalan}



\begin{quote}
Before the new one: say the Hindi for breath, then the Hindi for an itch.
\end{quote}

\subsection*{You'll want to know: \hi{जलन}}
\textbf{\hi{जलन}} — \emph{jalan} — \textquotedblleft{}a burning feeling\textquotedblright{}.

\begin{cousinweb}
\textbf{\hi{जलन}} (\emph{jalan}) is a tadbhava, a word inherited from Sanskrit and reshaped by sound change. It comes from \emph{jvalana}, \textquotedblleft{}burning\textquotedblright{} (CDIAL 5307).
\end{cousinweb}

\begin{grammarlens}[title={What you've built}]
One more word for this chapter.
\end{grammarlens}

\subsection*{Your turn}
\begin{itemize}
\raggedright
  \item \emph{Say it:} \emph{jalan}
  \item \emph{Say it:} \emph{jalan}, once more
  \item \emph{Say it:} say \emph{jalan}
  \item \emph{Recall:} say the Hindi for breath, then the Hindi for an itch, then say \emph{jalan} again
\end{itemize}

\begin{tcolorbox}[breakable,colback=teal!4,colframe=teal!35!black,title={Before you move on}]
What does \hi{जलन} mean? (\textquotedblleft{}A burning feeling\textquotedblright{}.) Say it once more.
\end{tcolorbox}
\section[\texorpdfstring{\hi{सूजन} (sūjan)}{सूजन (sūjan)}]{\hi{सूजन} (sūjan) — a swelling}

\label{lesson:HI-C274-sujan}



\begin{quote}
Before the new one: say the Hindi for an itch, then the Hindi for a burning feeling.
\end{quote}

\subsection*{You'll want to know: \hi{सूजन}}
\textbf{\hi{सूजन}} — \emph{sūjan} — \textquotedblleft{}a swelling\textquotedblright{}.

\begin{cousinweb}
\textbf{\hi{सूजन}} (\emph{sūjan}) is the noun of \emph{sūjnā}, \textquotedblleft{}to swell\textquotedblright{}, which comes from Sanskrit \emph{śūyate}, \textquotedblleft{}is swollen\textquotedblright{} (CDIAL 12568).
\end{cousinweb}

\begin{grammarlens}[title={What you've built}]
One more word for this chapter.
\end{grammarlens}

\subsection*{Your turn}
\begin{itemize}
\raggedright
  \item \emph{Say it:} \emph{sūjan}
  \item \emph{Say it:} \emph{sūjan}, once more
  \item \emph{Say it:} say \emph{sūjan}
  \item \emph{Recall:} say the Hindi for an itch, then the Hindi for a burning feeling, then say \emph{sūjan} again
\end{itemize}

\begin{tcolorbox}[breakable,colback=teal!4,colframe=teal!35!black,title={Before you move on}]
What does \hi{सूजन} mean? (\textquotedblleft{}A swelling\textquotedblright{}.) Say it once more.
\end{tcolorbox}
